\chapter{Estado del Arte}
En este apartado se estudiarán las soluciones existentes en 
teleoperación y visualización de 
parámetros de robots 
mediante computador, interfaces móviles y realidad aumentada.


\section{Soluciones de escritorio}
En cuanto a la operación telemática y visualización de parámetros 
de robots mediante computador existen 
distintas soluciones que ofrecen diferentes funcionalidades. Entre ellas 
cabe destacar \textit{RViz}, \textit{Foxglobe Studio} o \textit{Webviz}.

\textit{RViz} es el visualizador en tres dimensiones (3D) oficial 
de escritorio para el marco
de trabajo (\textit{framework}) ROS \parencite{rosRVizUser}.
La aplicación permite cargar un modelo del robot, nubes de puntos, mapas,
imágenes, estados del robot, etc. configurar diferentes cámaras en 
la visualización del gemelo del robot en el computador y realizar distintas 
acciones como seleccionar elementos en el panel de trabajo, centrar las 
cámaras a una localización especificada, medir la distancia existente
entre dos puntos del espacio, establecer el punto inicial del sistema 
de localización y el objetivo, etc. Además, permite realizar simulaciones 
de manera que los operadores puedan validar las acciones previamente
a ejecutarlas en el entorno real.

A pesar de que \textit{RViz} puede resultar una herramienta realmente útil
para la simulación o teleoperación de sistemas robóticos, la digitalización 
de los modelos 3D de los robots puede ser compleja para usuarios que 
no cuenten con conocimiento de modelado por computador o del formato 
URDF (\textit{Unified Robotics Description Format}), por lo 
que los resultados finales pueden no ser los más óptimos. Además, 
el entorno de operación de la aplicación está basado en la recreación del 
entorno real mediante gemelos digitales, lo que impide la operación 
directa sobre la robótica instalada y necesita de navegación 
por diferentes paneles y ventanas para modificar parámetros, lo que puede generar confusiones o 
excesivas cargas cognitivas para operadores no experimentados o 
familiarizados con la herramienta.

\textit{Foxglobe Studio} es una plataforma para registrar, subir, organizar
y visualizar datos de registro multimodales como series de tiempo, 
registros de texto, vídeo, 3D, mapas y otros \parencite{foxgloveFoxgloveDocumentation}.
La aplicación permite la integración con ROS2 mediante \textit{Foxglobe Bridge}. 
De manera similar al caso anterior, es posible cargar modelos, audio, mapas, registros 
de texto, etc. y además de otras operativas, permite la teleoperación básica publicando 
mensajes en un \textit{topic} específico.

A diferencia de \textit{RViz}, \textit{Foxglobe Studio} no solo permite 
utilizar otros \textit{framework} como \textit{PX4} o \textit{LeRobot}, sino 
que se trata de una aplicación más compleja con mayores funcionalidades que 
permite una gestión de grano fino de los diferentes procesos de la visualización 
de parámetros y gestión de datos. Sin embargo, se pueden encontrar las mismas problemáticas
que con la herramienta anterior, potenciando, en este caso, la complejidad de utilización
de la aplicación por estar más enfocada a procesos industriales en el que los 
operadores cuentan con mayor experiencia, de manera que puede ser inconveniente para 
usuarios no instruidos adecuadamente.

\textit{Webviz} permite visualizar parámetros del robot empleando el navegador \parencite{webvizWebviz},
lo que permite cruzar la frontera de los dispositivos de mayor potencia como 
portátiles y computadores de escritorio hacia dispositivos móviles como \textit{smartphones}
y \textit{tablets}. Además, se trata de una herramienta modular, de modo que se pueden 
añadir nuevos componentes que permitan visualizar otros datos diferentes.

A pesar de que esta aplicación permite ser empleada en todos los dispositivos al 
tratarse de una herramienta web, cuenta con problemáticas similares a las mencionadas
anteriormente.


\section{Soluciones móviles}
En cuanto a la operación telemática y visualización de parámetros empleando 
dispositivos móviles se pueden encontrar soluciones como \textit{ROS-Mobile}, 
\textit{Robot Steering App}, \textit{ROS2 Connect Mobile} o soluciones 
propietarias de ciertos fabricantes de robots.

\textit{ROS-Mobile} es una aplicación Android diseñada para el control dinámico y 
visualización de sistémas robóticos móviles ROS \parencite{rosROSMobileWiki} basada
en menús para la visualización de parámetros y la operación del robot.

\textit{Robot Steering App} es una aplicación iOS que permite controlar 
la velocidad lineal y angular, la gestión de los \textit{topics}, ejecutar paradas 
de emergencia y obtener flujos de vídeo en tiempo real de las cámaras de los robots ROS2
\parencite{appleRobotSteering}.

\textit{ROS2 Connect Mobile} proporciona una interfaz móvil para el monitoreo 
de \textit{topics} de ROS2 \parencite{appleROS2Connect}. A diferencia de las otras aplicaciones
no permite teleoperar los robots.

Las aplicaciones existentes en su mayoría se tratan de soluciones básicas 
para la visualización y operación en entornos limitadas en 
la interacción espacial empleando, además, 
navegación mediante menús, lo que aumenta la carga cognitiva del operador.


\section{Soluciones con realidad aumentada}
Existen diferentes soluciones que emplean realidad aumentada para la visualización de 
parámetros y teleoperación de robots. Sin embargo, en su mayoría se tratan de herramientas
diseñadas para solucionar un problema concreto empleando 
dispositivos HMD (\textit{Head Mounted Display}) en lugar de tratarse de aplicaciones 
multidominio que puedan ser aplicadas en diferentes entornos y dispositivos.

\cite{adekoya2025horusmixedrealityinterface} introduce HORUS, una interfaz de realidad aumentada 
para gestionar equipos de robots móviles empleando \textit{Meta Quest 2}.

\cite{10.1145/3610978.3640571} introduce AR-STAR 
una herramienta de realidad aumentada
que permite modificar nubes de puntos empleando sobre superficies
el visor \textit{HoloLens 2} para la reparación
de superficies corrodidas mediante robots tele-operados.

\cite{alietal} proponen una solución basada en realidad aumentada para la 
operativa telemática y visualización de datos de robots quirúrjicos empleando
\textit{HoloLens 2} para controlar el movimiento de la cabeza, las 
órdenes verbales y proporcionar la retroalimentación al usuario.

\cite{nevesetal} introduce una implementación de interfaz de realidad mixta mediante 
\textit{Microsoft HoloLens} con para proporcionar 
información sobre la tarea que se está realizando e interactuando con 
objetos virtuales como si se estuviera interactuando con objetos reales.

\cite{10.1145/3708468.3715687} propone una demostración de colaboración robot-humano empleando 
realidad extendida empleando \textit{Meta Quest 3}.

\cite{inbook} introduce una solución para integrar los robots industriales,
la programación industrial \textit{offline} y herramientas 
de simulación, plataformas de optimización y una interfaz de realidad aumentada
empleando \textit{HoloLens} de manera que se permita al operador ser una parte
más del sistema de producción robotizado.

\cite{10770031} desarrollan una solución para el control de entornos de construcción
mediante sensores y camaras integrados en conos delimitadores y una interfaz 
de realidad virtual para HMD que permite situar al usuario en el terreno.

Los terminales móviles como \textit{smartphones} y \textit{tablets} 
también han sido empleados, en menor medida, como interfaz humano-robot.

\cite{friske2024integrationaugmentedrealitymobile} propone una solución para proporcionar información del entorno contextualizada
en situaciones de emergencia mediante la comunicación con el robot y la visualización
de información en la pantalla del terminal.

\cite{10.22260/ISARC2019/0084} emplea realidad mixta para proporcionar visualización de 
datos gráficos de modelado de información de construcción
(BIM; \textit{Building Information Modeling}) en terminal móvil.

\cite{ROSALES2021618} introduce una solución de realidad aumentada para 
la recolección y visualización de datos en tiempo real 
en terminal móvil empleando sensores 
del internet de las cosas (IoT) en una fábrica de la automovilística Volvo.

Como puede apreciarse, la mayoría de las soluciones móviles que emplean realidad 
aumentada en dispositivos móviles están centradas en resolver un problema específico. 
\cite{ZEA2021100057}, sin embargo, introduce una aplicación para la visualización de datos y 
teleoperación en ROS1
empleando dispositivos móviles (Android/iOS/\textit{Microsoft HoloLens}), lo que 
permite visualizar los diferentes datos y parámetros de una manera más directa. 


